\begin{lemma}
Assume
\(\noderef{ParameterHierarchyEventualComparisons}
(\varepsilon,\varepsilon_1,\varepsilon_2,n_0)\), \(n>n_0\),
\[
k=\noderef{TwoBiteNaturalIndependenceScale}(1+\varepsilon,n),
\qquad |I|=k,
\]
and \(\noderef{FiberAndDegreeEvent}\).  Let \(H_I\) be the huge class and put
\[
t_1=\noderef{TwoBiteHugeCutoff}(n),\qquad C=100(\log n)^3.
\]
Then
\[
\sum_{x\in H_I}|X_x(I)|\le 2k,
\qquad
|H_I|\le \frac{2k}{t_1}.
\]
\end{lemma}
\begin{proof}
All finite cardinalities in this proof are used in two forms.  Before casting,
they are natural-number cardinalities of finite sets; when they occur in an
inequality with \(k,t_1\), or \(C\), they are cast to \(\mathbb R\).  The
identity \(|I|=k\) is therefore used as the real identity
\((|I|:\mathbb R)=(k:\mathbb R)\).

Recall that \(H_I\) is the finite filter of base vertices satisfying
\(\noderef{TwoBiteHugeClass}(\omega,I,x)\).  Thus, if \(x\in H_I\), the
membership statement first gives the predicate
\(\noderef{TwoBiteHugeClass}(\omega,I,x)\), and unfolding that predicate gives
the strict lower cutoff and the trivial upper bound
\[
t_1<|X_x(I)|\le |I|=k. \tag{1}
\]
Unfolding \(\noderef{TwoBiteX}\) gives the actual set identity
\[
X_x(I)=I\cap N_x,
\]
so every element of \(X_x(I)\) lies in \(I\).  Hence
\[
X_x(I)\subseteq I. \tag{1a}
\]

If \(x,y\in H_I\) are distinct, then (1a) is not needed for the overlap
bound; instead the definition \(X_z(I)=I\cap N_z\) gives
\[
X_x(I)\cap X_y(I)\subseteq N_x\cap N_y.
\]
We check the possible colors of \(x\) and \(y\), writing base vertices as
elements of the disjoint sum \(V_R\sqcup V_B\).  If
\(x=\operatorname{inl}(r)\) and \(y=\operatorname{inl}(s)\) are both red, then
\(x\ne y\) is exactly \(r\ne s\), and the red-red lifted-neighborhood
intersection clause of \(\noderef{FiberAndDegreeEvent}\) gives
\[
|N_r\cap N_s|\le 100(\log n)^3=C.
\]
If \(x=\operatorname{inr}(b)\) and \(y=\operatorname{inr}(c)\) are both blue,
the blue-blue clause gives the same bound for \(|N_b\cap N_c|\).  If one
vertex is red and the other blue, then the sum constructors already make them
distinct, and the mixed red-blue lifted-neighborhood intersection clause gives
the same bound.  In all three cases, cardinality monotonicity applied to
\(X_x(I)\cap X_y(I)\subseteq N_x\cap N_y\) gives
\[
|X_x(I)\cap X_y(I)|\le C. \tag{2}
\]
Thus, for every finite \(B\subseteq H_I\),
\[
k\ge\left|\bigcup_{x\in B}X_x(I)\right|
\ge \sum_{x\in B}|X_x(I)|
  -\binom{|B|}{2}_{\mathbb R}C. \tag{3}
\]
We spell out the second inequality.  Enumerate \(B=\{b_1,\ldots,b_r\}\) and
write \(A_i=X_{b_i}(I)\).  If \(r=0\), then the union and the sum are both
empty and the claimed lower bound is \(0\ge0\).  Assume now \(r>0\), and put
\[
U_j=A_1\cup\cdots\cup A_j\qquad (0\le j\le r),
\]
with \(U_0=\varnothing\).  For \(0\le j<r\), the finite-cardinality identity
\[
|U_{j+1}|=|U_j|+|A_{j+1}|-|U_j\cap A_{j+1}|.
\]
holds before casting.  After casting to real numbers we use the same equality
in \(\mathbb R\).  Since
\[
U_j\cap A_{j+1}\subseteq
\bigcup_{1\le i\le j}(A_i\cap A_{j+1}),
\]
the finite union bound gives
\[
|U_j\cap A_{j+1}|
\le\sum_{1\le i\le j}|A_i\cap A_{j+1}|.
\]
The elements \(b_i\) are all distinct, so (2) applies to each pair
\((b_i,b_{j+1})\).  Therefore, after casting all cardinalities and \(j\) to
\(\mathbb R\),
\[
|U_j\cap A_{j+1}|
\le jC.
\]
Inducting on \(j\) gives
\[
\left|\bigcup_{x\in B}X_x(I)\right|
\ge
\sum_{x\in B}|X_x(I)|
-
\sum_{j=0}^{r-1} jC
=
\sum_{x\in B}|X_x(I)|
-\binom r2_{\mathbb R}C.
\]
The last equality is the ordinary finite pair count:
\[
\sum_{j=0}^{r-1}j=\frac{r(r-1)}2=\binom r2_{\mathbb R},
\]
which is the number of unordered pairs in the \(r\)-element enumeration.
Here \(r=|B|\), so this is exactly the lower bound in (3).  The first
inequality in (3) follows because (1a) gives \(X_x(I)\subseteq I\) for every
\(x\), so the union over \(B\) is a subset of \(I\), and \(|I|=k\).

Apply the eventual-comparisons package at the present \(n\).  Its
huge-overlap conjunct is the explicit finite inequality
\[
\binom{2k/t_1+1}{2}_{\mathbb R}C\le\frac12k. \tag{4}
\]
The cutoff \(t_1\) is positive for the present large \(n\): unfolding
\(\noderef{TwoBiteHugeCutoff}\) gives
\[
t_1=\frac{\sqrt{n\log n}}{\log\log n},
\]
and the eventual range is chosen beyond \(e^e\), so \(n>1\),
\(\log n>0\), \(\log\log n>0\), and the quotient is positive.  Also
\(C=100(\log n)^3\) is nonnegative, and \(k=(|I|:\mathbb R)\ge0\).  Hence
\[
a:=2k/t_1
\]
is nonnegative.  Suppose first that
\[
(|H_I|:\mathbb R)>a. \tag{5}
\]
Let \(r=\lceil a\rceil\), the least natural number with \(a\le r\).  Since
\(|H_I|\) is a natural number and its real cast is strictly larger than \(a\),
we have
\[
r\le |H_I|. \tag{6}
\]
Also the ceiling property gives
\[
a\le r\le a+1. \tag{7}
\]
If \(r=0\), then \(a=0\), and (5) says \(H_I\) is nonempty.  Choose
\(x\in H_I\).  Since \(a=2k/t_1\), \(a=0\), \(t_1>0\), and \(k\ge0\), we have
\(k=0\).  Then (1) gives \(t_1<|X_x(I)|\le k=0\), contradicting \(t_1>0\).
Thus \(r>0\).

By (6), choose a subfamily as follows.  List the finite set \(H_I\) without
repetition as \(h_1,\ldots,h_q\), where \(q=|H_I|\).  Since \(r\le q\), take
\[
B=\{h_1,\ldots,h_r\}.
\]
Then
\[
B\subseteq H_I,\qquad |B|=r=\left\lceil 2k/t_1\right\rceil.
\]
For each \(x\in B\), (1) gives the strict real inequality
\(|X_x(I)|>t_1\).  Because \(B\) has \(r>0\) elements, summing these strict
inequalities is legitimate: the sum of the positive gaps
\(\sum_{x\in B}(|X_x(I)|-t_1)\) is positive.  Hence
\[
\sum_{x\in B}|X_x(I)|>rt_1. \tag{8}
\]
Since \(t_1>0\) and \(r\ge a=2k/t_1\), we have
\[
rt_1\ge 2k. \tag{9}
\]
Moreover \(r\le a+1\).  The real binomial function
\(\binom u2_{\mathbb R}=u(u-1)/2\) is increasing on \([1,\infty)\): if
\(1\le u\le v\), then
\[
\binom v2_{\mathbb R}-\binom u2_{\mathbb R}
  =\frac{(v-u)(u+v-1)}2\ge0.
\]
If \(r=1\), then \(\binom r2_{\mathbb R}=0\), while
\(a+1\ge1\) and \(C\ge0\) make
\(\binom{a+1}{2}_{\mathbb R}C\ge0\).  Therefore, using \(C\ge0\),
\[
\binom r2_{\mathbb R}C
\le
\binom{a+1}{2}_{\mathbb R}C
\le \frac12 k, \tag{10}
\]
where in the case \(r\ge2\) the first inequality uses the displayed
monotonicity with \(1\le r\le a+1\), and the last inequality is exactly (4).
Combining (3), (8), (9), and (10),
\[
k\ge\left|\bigcup_{x\in B}X_x(I)\right|
> |B|t_1-\binom{|B|}{2}_{\mathbb R}C
\ge 2k-\frac12k=\frac32k,
\]
a contradiction because \(k\ge0\) implies \(\frac32k\ge k\).  Hence
\[
(|H_I|:\mathbb R)\le 2k/t_1. \tag{11}
\]
This is the second asserted bound.

Now take \(B=H_I\) in (3).  Since
\((|H_I|:\mathbb R)\le2k/t_1=a\), the same comparison (4), with the same
monotonicity justification as above, gives the needed overlap estimate.
Indeed, if \(|H_I|=0\) or \(|H_I|=1\), then
\(\binom{|H_I|}{2}_{\mathbb R}=0\), and hence
\[
\binom{|H_I|}{2}_{\mathbb R}C=0\le\frac12k
\]
because \(C\ge0\) and \(k\ge0\).  Otherwise \(|H_I|\ge2\), so
\(|H_I|\) lies in the monotonicity interval \([1,\infty)\).  From (11),
\(|H_I|\le a\le a+1\), and the monotonicity of
\(\binom u2_{\mathbb R}\) on \([1,\infty)\), followed by multiplication by
the nonnegative number \(C\) and then by (4), gives
\[
\binom{|H_I|}{2}_{\mathbb R}C\le\frac12k.
\]
The lower bound (3) for \(B=H_I\) rearranges to
\[
\sum_{x\in H_I}|X_x(I)|
\le k+\binom{|H_I|}{2}_{\mathbb R}C.
\]
Therefore
\[
\sum_{x\in H_I}|X_x(I)|
\le k+\binom{|H_I|}{2}_{\mathbb R}C
\le \frac32k\le2k.
\]
This proves both asserted bounds.
\end{proof}
